% =====================================================================
\chapter{Diventare autonomi}
% =====================================================================

Questa appendice non parla di C++: parla di \textbf{come studiare e lavorare} da solo. In
laboratorio, all'esame, sarai tu, un editor, un compilatore e una traccia: niente
Internet, niente AI, niente dispensa. Le abitudini che costruisci in queste settimane
decidono se in quel momento ti sentirai perso o a casa. E sono le stesse che, nel lungo
periodo, separano chi \emph{usa} il codice degli altri da chi sa \emph{scriverlo}.

% ---------------------------------------------------------------------
\section{Come si studia la programmazione}
% ---------------------------------------------------------------------

\begin{analogia}[: la palestra]
Guardare qualcuno che solleva pesi non ti fa venire i muscoli. Leggere codice scritto
bene, o una spiegazione chiara (anche questa dispensa), ti dà la sensazione di aver
capito: è la sensazione di chi guarda la palestra dalla vetrina. Il muscolo si forma
solo quando \textbf{sei tu} a sollevare: a scrivere il codice da zero, a sbagliare, a
capire perché.
\end{analogia}

\subsection*{Le cinque tecniche che funzionano davvero}

\begin{passopasso}[: le tecniche]
\begin{enumerate}
  \item \textbf{Richiamo attivo (la pagina bianca).} Dopo aver studiato un argomento,
        chiudi tutto e scrivi su un foglio bianco quello che ricordi: definizioni, un
        disegno della memoria, il codice. Poi confronta. Quello che non sei riuscito a
        scrivere è quello che non sai ancora. Rileggere sottolineando, al contrario, dà
        l'illusione di sapere.
  \item \textbf{Riscrivere da zero.} Il pattern di ogni modulo va scritto \emph{senza
        guardarlo} finché non esce. La prima volta servono 40 minuti, la terza 10. A
        quel punto in laboratorio la lista, la pila o il BST escono \q{dalle dita}, e la
        testa resta libera per la parte nuova della traccia.
  \item \textbf{Ripetizione distanziata.} Si dimentica in fretta ciò che si è visto una
        volta sola. Ripassa un argomento dopo 1 giorno, poi dopo 3, poi dopo 7: per
        questo il piano di studio ha ogni giorno 15 minuti di quiz dei moduli passati.
  \item \textbf{Spiegare ad alta voce (tecnica di Feynman).} Spiega l'argomento come se
        lo dovessi insegnare a un amico che non sa programmare. Dove ti incagli, lì c'è un
        buco. È anche il miglior allenamento per l'orale.
  \item \textbf{Mescolare.} Negli ultimi giorni non fare un modulo alla volta: alterna
        quiz di argomenti diversi. All'esame le domande arrivano mescolate, e devi
        riconoscere \emph{quale} strumento serve, non solo saperlo usare.
\end{enumerate}
\end{passopasso}

\subsection*{Carta e penna prima della tastiera}

Prima di scrivere una funzione su liste, alberi o grafi, \textbf{disegna}: i nodi, i
puntatori, cosa cambia passo per passo. Quasi tutti gli errori sulle strutture dinamiche
sono puntatori aggiornati nell'ordine sbagliato, e su un disegno si vedono subito. Le
figure TikZ di questa dispensa sono il modello: rifalle a mano.

\subsection*{Studiare con poco tempo e molta stanchezza}

\begin{itemize}
  \item \textbf{Blocchi da 50 minuti} con 10 minuti di pausa vera (in piedi, lontano dallo
        schermo). Dopo 4 blocchi, una pausa lunga. La concentrazione non dura 8 ore di
        fila per nessuno.
  \item \textbf{Il compito più difficile al mattino}: codice e tracce quando sei fresco;
        quiz e ripasso quando sei stanco.
  \item \textbf{Il sonno fa parte dello studio}: è di notte che la memoria si consolida.
        Una notte in bianco prima dell'esame costa più di quanto rende.
  \item \textbf{Telefono in un'altra stanza}. Ogni notifica costa minuti di concentrazione.
\end{itemize}

% ---------------------------------------------------------------------
\section{Usare l'AI il meno possibile (e bene)}
% ---------------------------------------------------------------------

L'AI è uno strumento potentissimo, e proprio per questo è pericolosa mentre si
\emph{impara}: risolve il problema al posto tuo, e ti toglie esattamente lo sforzo che ti
avrebbe insegnato qualcosa. All'esame non ci sarà. L'obiettivo di queste settimane è
arrivare a non averne bisogno.

\begin{analogia}[: il navigatore]
Chi guida sempre con il navigatore, dopo anni non sa ancora come arrivare in centro. Chi
ha girato qualche volta con la cartina e si è perso, la strada la conosce. Usa l'AI come
la cartina che consulti quando sei davvero perso, non come la voce che ti dice
\q{tra 200 metri gira a destra} a ogni incrocio.
\end{analogia}

\subsection*{La scala dell'autonomia: cosa fare quando sei bloccato}

Prima di chiedere aiuto a chiunque (AI, un collega, la soluzione), sali questi gradini
nell'ordine. Ognuno risolve da solo la maggior parte dei problemi.
\begin{passopasso}[: quando sei bloccato]
\begin{enumerate}
  \item \textbf{Leggi l'errore del compilatore}, dall'\emph{alto}: il primo errore è
        quello vero, gli altri spesso sono conseguenze. Guarda il numero di riga, e
        anche la riga prima (il punto e virgola mancante si segnala sulla riga dopo).
  \item \textbf{Riduci il problema}: commenta metà del codice. Se l'errore sparisce, è
        nell'altra metà. Ripeti: è la ricerca binaria applicata ai bug (Modulo 7).
  \item \textbf{Stampa}: un \cc|cout| prima e dopo il punto sospetto, con i valori delle
        variabili e dei puntatori. Quasi ogni crash su liste e alberi è un
        \cc|nullptr| dereferenziato: stampa i puntatori.
  \item \textbf{Disegna} lo stato della memoria a quel punto, con carta e penna.
  \item \textbf{Spiega il codice riga per riga ad alta voce}, anche a un oggetto (è la
        famosa \q{tecnica della paperella di gomma}): a metà spiegazione, spesso, vedi
        l'errore.
  \item \textbf{Torna al modulo} della dispensa: il pattern di codice e le trappole.
  \item \textbf{La regola dei 20 minuti}: se dopo 20 minuti di tentativi seri sei ancora
        fermo, \emph{allora} puoi chiedere aiuto. Non prima (non hai provato abbastanza),
        non molto dopo (stai solo perdendo tempo).
\end{enumerate}
\end{passopasso}

\subsection*{Se usi l'AI, usala come un tutor, non come uno che fa i compiti}

\begin{itemize}
  \item Chiedi \textbf{un indizio}, non la soluzione: \q{non dirmi il codice: dove devo
        guardare?}, \q{che cosa non ho capito di questo concetto?}.
  \item Chiedi di \textbf{spiegarti l'errore}, non di correggerlo: \q{cosa significa
        questo messaggio del compilatore?}.
  \item Dopo aver ricevuto un aiuto, \textbf{chiudi la conversazione e riscrivi tutto da
        solo}, senza guardare. Se non ci riesci, non l'hai ancora capito.
  \item Usala per \textbf{interrogarti}: \q{fammi 5 domande d'orale sulla Regola dei Tre e
        correggi le mie risposte}. Qui l'AI è un ottimo allenatore.
  \item \textbf{Mai} copiare codice che non sapresti riscrivere e spiegare riga per
        riga. All'orale ti chiederanno proprio quello.
\end{itemize}

\begin{trappola}[: l'illusione di competenza]
Leggere una soluzione corretta (dell'AI o di una dispensa) produce la sensazione
\q{sì, l'avrei fatto anch'io}. Quasi sempre è falsa. L'unica prova è riscriverla da zero,
il giorno dopo, senza guardare.
\end{trappola}

\subsection*{Gli strumenti che rendono autonomi}

\begin{itemize}
  \item \textbf{Il compilatore con gli avvisi}: \texttt{g++ -Wall -Wextra}. Gli avvisi
        sono errori futuri: variabili non inizializzate, confronti sbagliati, \cc|return|
        mancanti. Tutto il codice di questa dispensa compila senza avvisi.
  \item \textbf{Il debugger} (\texttt{gdb}, o quello dell'IDE): fermare il programma su una
        riga, guardare le variabili, avanzare un'istruzione alla volta. Imparalo: 30
        minuti spesi oggi risparmiano ore. Compila con \texttt{-g}.
  \item \textbf{La documentazione}: \texttt{cppreference.com} per il C++. Leggere la
        documentazione è un'abilità che si allena come le altre.
  \item \textbf{Il tuo quaderno degli errori}: ogni volta che un bug ti fa perdere più di
        10 minuti, scrivi in una riga cos'era e come l'hai trovato. Rileggilo prima
        dell'esame: i tuoi errori tendono a ripetersi.
\end{itemize}

% ---------------------------------------------------------------------
\section{Superare l'esame}
% ---------------------------------------------------------------------

Strategie per ogni prova, piano minuto per minuto del laboratorio, checklist
prima della consegna, aiuti disponibili e cosa fare se un appello va male sono
raccolti nel capitolo \q{Superare l'esame: tutte le strade}, subito prima di
questo. Qui resta solo il principio che li riassume tutti: \textbf{allenati nelle
stesse condizioni dell'esame}, da solo, con il cronometro, senza Internet e senza AI.

% ---------------------------------------------------------------------
\section{Diventare un buon programmatore (nel lungo periodo)}
% ---------------------------------------------------------------------

L'esame è una tappa. Questi sono i consigli che valgono per i prossimi anni.

\subsection*{Le abitudini}

\begin{enumerate}
  \item \textbf{Scrivi codice ogni giorno}, anche poco. Venti minuti al giorno battono
        cinque ore la domenica. La programmazione è un'abilità, come uno strumento
        musicale: si mantiene solo usandola.
  \item \textbf{Costruisci progetti tuoi}. Gli esercizi insegnano gli strumenti; i progetti
        (un gestionale per una cosa che ti interessa, un piccolo gioco, uno script che
        automatizza un tuo lavoro noioso) insegnano a \emph{usarli insieme} e a decidere da
        soli. È lì che si diventa autonomi davvero.
  \item \textbf{Leggi codice scritto da altri}, bene: progetti open source, il codice dei
        colleghi. Chiediti perché è scritto così. Si impara a scrivere leggendo.
  \item \textbf{Scegli bene i nomi}. \cc|contaStudentiPromossi| invece di \cc|f2|.
        Il codice si legge molte più volte di quante si scriva: chi lo leggerà tra sei
        mesi sarai tu.
  \item \textbf{Funzioni corte che fanno una cosa sola}. Se per descrivere una funzione
        ti serve la parola \q{e}, probabilmente sono due funzioni.
  \item \textbf{Prova il tuo codice}: pensa sempre ai casi limite (vuoto, uno solo,
        duplicati, valori enormi) e scrivi piccoli test. Il codice non provato è codice
        rotto che non lo sa ancora.
  \item \textbf{Usa Git} fin da subito, anche per i progetti personali: salvi la storia
        del tuo lavoro, puoi sperimentare senza paura e torni indietro quando rompi
        qualcosa. È uno strumento che ti chiederanno ovunque.
  \item \textbf{Prima fallo funzionare, poi rendilo bello, poi (se serve) veloce}. In
        quest'ordine.
\end{enumerate}

\subsection*{Cosa studiare dopo}

\begin{itemize}
  \item \textbf{Algoritmi e strutture dati}: sono il fondamento di tutto, e questo corso
        ne è l'inizio. Allenati con problemi di programmazione (siti di esercizi come
        quelli delle gare di programmazione, o l'\emph{Advent of Code} ogni dicembre).
  \item \textbf{Il C++ moderno}: smart pointer (\cc|unique_ptr|, \cc|shared_ptr|), la STL
        che qui era vietata (\cc|vector|, \cc|map|), RAII. Ora che sai cosa c'è sotto, li
        userai capendo cosa fanno.
  \item \textbf{Un secondo linguaggio}, diverso (ad esempio Python): vedere le stesse idee
        in un'altra forma le fa capire meglio. Ma uno alla volta, e solo dopo aver
        consolidato il primo.
  \item \textbf{Come funziona il computer}: sistemi operativi, reti, basi di dati. Sono i
        corsi che verranno, e danno senso a tutto il resto.
\end{itemize}

\subsection*{Il rapporto con l'AI, dopo}

Quando le basi sono solide, l'AI diventa un moltiplicatore: ti fa risparmiare il lavoro
ripetitivo, e tu sai giudicare se quello che produce è giusto. Senza basi, invece, ti
rende dipendente da qualcosa che non sai controllare. La regola resta la stessa:
\textbf{non usare mai codice che non sapresti scrivere e spiegare da solo}. Chi capisce
davvero il codice sarà sempre più prezioso di chi sa solo chiederlo.

\begin{inparole}
\begin{itemize}
  \item Si impara scrivendo, non leggendo: pagina bianca, riscrivere da zero, spiegare
        ad alta voce.
  \item Quando sei bloccato: errore del compilatore, ridurre, stampare, disegnare,
        spiegare, dispensa; solo dopo 20 minuti chiedi aiuto.
  \item AI come tutor che ti interroga e ti dà indizi, mai come chi fa i compiti al posto
        tuo.
  \item Per l'esame: gerarchia + strutture dati a memoria, 5--6 tracce vere a tempo,
        orale ad alta voce, sonno.
  \item Per la carriera: codice ogni giorno, progetti tuoi, nomi chiari, test, Git.
\end{itemize}
\end{inparole}

\vspace{1em}
\begin{center}
\large\itshape In bocca al lupo per l'esame!
\end{center}
